% TOP_PROBLEM: {"id": 2100514, "title": "Open Problems in Integrable Systems — Nijenhuis operators: singular points and global properties", "category_id": 11, "category": {"id": 11, "name": "dynamical_systems", "display_name": "Dynamical Systems", "description": "Problems about long-term behavior of deterministic systems, Hamiltonian dynamics, and periodic orbits.", "slug": "dynamical-systems", "order_index": 11, "created_at": "2026-07-31T15:26:25.671Z"}, "status": "open"}
% FIRST_POSTED: null
% ENTRY_KIND: statement_only
% Imported from: batch09.tex; UnsolvedMath ID 2100514
\documentclass[11pt]{article}
\usepackage[margin=25mm]{geometry}
\usepackage{fontspec}
\setmainfont{Latin Modern Roman}
\usepackage{amsmath,amssymb,mathtools}
\usepackage{unicode-math}
\setmathfont{Latin Modern Math}
\usepackage{xurl,hyperref}
\hypersetup{colorlinks=true,urlcolor=blue}
\setcounter{secnumdepth}{0}
\setlength{\parindent}{0pt}
\setlength{\parskip}{5pt}
\setlength{\emergencystretch}{3em}
\newcommand{\sref}[2]{\hyperlink{source:#1:#2}{[S#2]}}
\newcommand{\cataloguescope}[1]{\paragraph{Recorded target.}#1}
\begin{document}
% UnsolvedMath ID 2100514; source catalogue code AMR-020-0514
\setcounter{section}{2100513}
\section{Open Problems in Integrable Systems — Nijenhuis operators: singular points and global properties}\label{2100514}
{\small\sffamily UnsolvedMath ID 2100514 · Dynamical Systems}
\subsection{Definitions and mathematical statement}

\paragraph{Shared notation.}$\mathbb N=\{1,2,\ldots\}$, and $\mathbb Z,\mathbb Q,\mathbb R,\mathbb C$ denote the integers, rationals, reals and complex numbers. Any other conventions are specified locally.


Let $(M^{2n},\omega)$ be symplectic and $P$ a Poisson tensor compatible with $\omega^{-1}$, meaning every linear combination of these two tensors is Poisson. The recursion operator is $R=P^{\sharp}\circ\omega^{\flat}:TM\to TM$ and is Nijenhuis. Take $n$ independent characteristic-coefficient functions $\sigma_1,\ldots,\sigma_n$ of $R$, assumed to have independent differentials almost everywhere, and put $\Phi=(\sigma_1,\ldots,\sigma_n)$. Generically $R$ is semisimple and every eigenvalue has multiplicity two. A singular point of $R$ is one where its Jordan type is not locally constant; a singular point of $\Phi$ has $\operatorname{rank}d\Phi<n$.\par A left-symmetric algebra is an algebra whose associator satisfies
\[a\cdot(b\cdot c)-(a\cdot b)\cdot c=b\cdot(a\cdot c)-(b\cdot a)\cdot c.\]
For a rank-$k$ critical point of an integrable system, Eliasson nondegeneracy means that the quadratic Hamiltonians induced by the critical components on the $2(n-k)$-dimensional symplectic normal space span a Cartan subalgebra of its symplectic Lie algebra.
\paragraph{Statement recorded in the supplied source.}
Describe the singularities of $\Phi$ in terms of the singular points of $R$, their linearizations and the associated left-symmetric algebras. Find sufficient and/or necessary conditions, in these terms, for a singularity of $\Phi$ to be Eliasson nondegenerate. Retain the special generic double-eigenvalue type of $R$. \sref{2100514}{2}
\subsection{Short English statement}
Relate singularities of the characteristic-coefficient integrable system to the recursion tensor and its linear algebra, and characterize when those singularities are nondegenerate.
\subsection{Sources}
\begin{description}
\item[\hypertarget{source:2100514:1}{S1}] ULAM.AI and the UnsolvedMath Contributors, UnsolvedMath: A Curated Collection of Open Mathematics Problems, record 2100514 (AMR-020-0514). CC BY 4.0. \url{https://huggingface.co/datasets/ulamai/UnsolvedMath}.
\item[\hypertarget{source:2100514:2}{S2}] Alexey Bolsinov, Vladimir S. Matveev, Eva Miranda and Serge Tabachnikov, Open Problems, Questions, and Challenges in Finite-Dimensional Integrable Systems (2018), Problem 5.17 and its complete preceding definitions. Original question and full discussion. \url{https://arxiv.org/abs/1804.03737}.

\end{description}
\label{end:2100514}
\end{document}
